\section{Small outputs}\label{sec:small}

The constructive attack needs the network output to be small relative to the gradient. There are two
routes to this, and the paper uses both because they cover different regimes. The first is purely
probabilistic and assumes nothing about the symmetry of the network or the input law: a random linear
readout makes the output small at every fixed input, and uniformly over inputs in the wide regime.
The second is the balance-and-concentration mechanism, which needs the classifier's sign to be
balanced but then applies to arbitrary --- in particular trained --- classifiers on the sphere. We
state both; the fragility theorem is proved once and fed by either.

\subsection{Random readout}
\begin{proposition}[Small output from a random readout]\label{prop:readout}
Let $F(x)=\langle w,h_L(x)\rangle$ with $w\sim N(0,I_{n_L}/n_L)$ independent of the other weights. For
each fixed $x$ and $\delta\in(0,1)$, with probability at least $1-\delta$ over $w$,
\[
  |F(x)|\ \le\ \|h_L(x)\|\,\sqrt{\tfrac{2}{n_L}\log\tfrac2\delta}.
\]
In the wide regime $d\log\bar n=o(n_L)$, w.h.p.\ over all the weights this holds simultaneously for
every $x$ on the sphere of radius $R$ with $\delta$ replaced by an exponentially small quantity, so
that
\[
  |F(x)|\ \le\ C\,\|h_L(x)\|\,\sqrt{\tfrac{d\log \bar n}{n_L}}\ \le\ C'\,\|x\|\,\sqrt{\tfrac{d\log\bar n}{n_L}}
  \qquad\text{for all }x .
\]
\end{proposition}

\begin{proof}
Conditionally on $h_L(x)$, $F(x)=\langle w,h_L(x)\rangle\sim N(0,\|h_L(x)\|^2/n_L)$, and a Gaussian tail
gives $\PP(|F(x)|>\|h_L\|\sqrt{2\log(2/\delta)/n_L})\le\delta$, which is the fixed-$x$ statement. For the
uniform statement, take an $\eps$-net $\mathcal N$ of the input sphere of radius $R$, of size
$|\mathcal N|\le(3R/\eps)^d$, and choose the threshold $\tau=\|h_L\|_{\max}\sqrt{2(d\log(3R/\eps)+1)/n_L}$
where $\|h_L\|_{\max}=\sup_x\|h_L(x)\|\le C\|x\|$ by the upper half of \eqref{eq:layerbracket}. The union
bound gives $\PP(\exists\,x_j\in\mathcal N:|F(x_j)|>\tau)\le(3R/\eps)^d\cdot2e^{-\tau^2n_L/(2\|h_L\|_{\max}^2)}\le2e^{-1}$,
and iterating the argument with $d\log(3R/\eps)$ replaced by $d\log(3R/\eps)+s$ makes the failure
probability $e^{-s}$, so w.h.p.\ $|F(x_j)|\le\tau$ for all net points. Finally $F$ is Lipschitz in $x$
with constant $\|w\|\prod_\ell\|W_\ell\|_\op=O(1)$ at fixed depth, so choosing $\eps=R/n$ extends the
bound from the net to all $x$ at the cost of $O(\|x\|/n)\ll\tau$, and $\log(3R/\eps)=\log(3n)=O(\log n)$.
Hence $|F(x)|\le C\|x\|\sqrt{(d\log n)/n_L}$ for all $x$ w.h.p. This is the union over the net at the
net-cardinality scale, not a single exponentially small $\delta$; the $\sqrt{d\log n}$ factor is the
price of the union.
\end{proof}

\Cref{prop:readout} is what removes the oddness and bias-free hypotheses of \cite{advExt}: no symmetry
of the network is used, the biases are arbitrary, and the activation is any $[0,1]$-slope function.
The price is that the readout must be random (or, for a trained network, close to a random draw in the
relevant inner products); the balance route below removes even that, at the cost of a symmetry
hypothesis on the input law.

\subsection{Balance and spherical concentration}
The second route is the geometric mechanism of \cite{shmaloAdv,advExt}, which we restate because the
fragility theorem uses it for the trained-network regime and because the certification results of
\Cref{sec:certificates} refer to it. Recall the Lévy constant $C_{\mathrm{L}}=\sqrt{\pi/8}$ and set,
for a $p$-balanced sign (meaning $\min(\mu(F>0),\mu(F<0))\ge p$),
\[
  r(p)=\frac{R}{\sqrt{d-1}}\sqrt{2\log(1/p)},\qquad
  \eps_{\mathrm{adv}}(p,\delta)=r(p)+\frac{R}{\sqrt{d-1}}\sqrt{2\log(2C_{\mathrm{L}}/\delta)}.
\]

\begin{proposition}[Small output from balance; {\cite[Thm 2.1]{advExt}}]\label{prop:balance}
Let $f:\Sph^{d-1}(R)\to\R$ be $L$-Lipschitz with $\sgn f$ $p$-balanced. Then for every $\delta\in(0,1)$,
on a set of $\mu$-measure at least $1-\delta$,
\[
  |f(x)|\ \le\ L\,\eps_{\mathrm{adv}}(p,\delta).
\]
If $f$ is built from bias-free linear maps, odd activations, and residual additions, then $f(-x)=-f(x)$
and the sign is exactly $\tfrac12$-balanced under any origin-symmetric input law with $\mu(f=0)=0$.
\end{proposition}

The two routes are complementary. The random-readout bound is pointwise (every $x$) and needs no
symmetry, but wants a random readout; the balance bound is measure-theoretic (all but a $\delta$
fraction) and needs a balanced sign, but then covers any Lipschitz classifier, including one obtained
by training, and any origin-symmetric input distribution. In the wide random regime both give
$|F(x)|=O(\|x\|/\sqrt{n})$ up to logarithms and a $\sqrt d$ factor for the uniform version. We use the
random-readout route for the clean every-$x$ statement in \Cref{sec:fragility} and note the balance
route as the trained-network version in \Cref{sec:trained}.

\begin{remark}[Why small output is the easy half]\label{rem:easy-half}
Neither route uses anything about the gradient; both are consequences of a Lipschitz function being
close to its median (concentration) and the median being near zero (balance, or a random readout with
mean zero). This is the observation of \cite{advExt} that the ``small output'' half of an adversarial
argument is generic, and that the architecture-specific work is entirely in the gradient lower bound
and its stability --- which is exactly what the submatrix theory of \Cref{sec:gains} provides.
\end{remark}
